\documentclass[11pt]{amsart}

\usepackage[T1]{fontenc}
\usepackage[utf8]{inputenc}
\usepackage{amsmath,amssymb,amsthm}
\usepackage[margin=1in]{geometry}
\usepackage{booktabs}
\usepackage{hyperref}

\newcommand{\LL}{\Lambda}
\newcommand{\OO}{\mathbb{O}}
\newcommand{\ZZ}{\mathbb{Z}}
\newcommand{\RR}{\mathbb{R}}
\newcommand{\Nrm}{N}
\newcommand{\Aut}{\mathrm{Aut}}
\newcommand{\Co}{\mathrm{Co}}

\theoremstyle{remark}
\newtheorem*{remarknote}{Remark}

\title[Update brief: Petersson reference; $S_3$ symmetry]%
{Update brief for\\
\emph{An order on the Leech lattice from a $\ZZ_3$-symmetric
triple-octonion product}\\
\large v4 (frozen 25 May 2026)}

\author{Jens K\"oplinger}
\date{26 May 2026}

\begin{document}
\maketitle

\begin{abstract}
The main paper is frozen at v4 (25 May 2026); the working source is
\texttt{paper/main.tex} and the frozen snapshot is
\texttt{paper/main\_2026-05-25.tex}.  This brief records two items
for a future revision: (1) the Petersson 2018 reference is presently
dangling in the bibliography and should be cited in the body text;
(2) the triple-octonion construction is in fact $S_3$-symmetric on
the three blocks, not merely $\ZZ_3$-symmetric.  Both items are
documented here against v4 and are not yet folded into the
manuscript.
\end{abstract}

% ===================================================================
\section{Petersson 2018 -- dangling reference}
% ===================================================================

\subsection*{Status}

The reference \texttt{Petersson2018} appears in the v4 bibliography
(H.\,P.~Petersson, ``Integral octonions,'' lecture at the M\'alaga
Workshop on Non-Associative Algebras, 2018) but is not cited
anywhere in the body of v4.  This is a regression from earlier
versions of the paper, where Petersson's lecture was cited as the
key modern account of the integral-octonion history.

\subsection*{Why this is important}

Petersson's 2018 lecture is the link in the historical chain that
brings the 1923--1946 narrative into the present:
\begin{itemize}
\item It is the source from which this project first identified
  Kirmse's $J_1$ as a specific non-closed candidate (and not the
  whole of Kirmse's work).
\item It treats Kirmse with appropriate care -- calling his work
  ``the bold attempt of a young mathematician'' -- and is therefore
  the reference for the fair-attribution stance the present paper
  takes in Appendix~B.
\item It also flags Dickson's 1923 priority, which the paper now
  records in Appendix~B's opening.
\end{itemize}

\subsection*{Proposed insertions for the next revision}

\paragraph{End of paragraph in \S2.2 (\emph{The $E_8$ lattice}).}
After the existing sentence ending ``\,closed under multiplication
and admits no enlargement to a larger order'' (citing Coxeter 1946
and Conway--Smith 2003), append a clause crediting Petersson:
\begin{quote}
``\ldots; for a modern survey of the integral-octonion history,
see Petersson (2018).''
\end{quote}

\paragraph{Appendix B.}
Petersson 2018 is the natural citation in the historical note where
the project's stance on Kirmse is established.  Two natural spots:
\begin{itemize}
\item In the \textbf{Kirmse (1924)} paragraph, after ``The remaining
  seven of his eight are the genuine maximal orders.'' add a
  parenthetical: ``A modern account of Kirmse's contribution,
  treating it as the bold work of a young mathematician on a topic
  weird and bizarre even today, is in Petersson (2018).''
\item Alternatively, at the start of Appendix B, in the opening
  sentence ``The construction this paper builds on was developed
  across four primary papers between 1923 and 1946,'' append a
  pointer to Petersson (2018) as the modern source from which this
  project first traced the chain.
\end{itemize}

Either placement (or both) closes the dangling reference and
restores the credit that earlier versions had.

% ===================================================================
\section{Kirmse 1924 -- two further contributions worth crediting}
% ===================================================================

The author's reading of Kirmse 1924 in the original highlights two
further contributions that the v4 Appendix B Kirmse paragraph does
not currently capture and that deserve credit.

\paragraph{(a) A nascent ideal theory on the $E_8$ lattice.}
Kirmse 1924 develops an ideal-theoretic treatment of the integral
octonions, working out divisibility and ideal-type structure on the
lattice he has constructed.  This is an early and substantive
companion contribution to the multiplicative-closure question that
the v4 Kirmse paragraph focuses on -- and it is the first such
development on an $E_8$-type lattice of integral octonions.  Mahler
1942 takes up the ideal-theoretic line again with the Euclidean
property and the principal-ideal theorem, but Kirmse's earlier
work in this direction is currently underrepresented.

\paragraph{(b) Algebra alternativity as a foundational building block.}
Kirmse takes \emph{alternativity} -- the weakened associativity
available in the octonions (and more generally, in alternative
algebras) -- as a basic axiom on which to build his treatment,
rather than working from full associativity and recovering
alternativity as a consequence.  This is a methodological choice
that deserves recognition in its own right: it suggests a
programme of replacing associativity by alternativity at the
foundation of an arithmetic of integral octonions, and it
anticipates the structural role alternativity plays in Cayley--
Dickson algebras and their integral arithmetic in the later
literature.

\subsection*{Proposed handling}

In a future revision, extend the v4 Appendix B Kirmse paragraph
with one or two sentences acknowledging both contributions.  A
draft:

\begin{quote}
``Kirmse also develops an ideal-theoretic treatment of the integral
octonions -- a substantive companion contribution that is taken up
again in Mahler (1942) -- and methodically takes algebra
alternativity as a foundational constraint, anticipating the
structural role alternativity plays in later integral-octonion
arithmetic.''
\end{quote}

This complements the Petersson 2018 citation proposed in \S1: the
Petersson reference establishes the project's fair-attribution
stance toward Kirmse, and these two sentences make that stance
concrete by naming the specific further contributions that the
narrowly-focused v4 paragraph leaves to the side.

% ===================================================================
\section{$S_3$ symmetry on the three blocks}
% ===================================================================

\subsection*{Finding}

The triple-octonion product $\star$ (Definition~3.3 of v4) is in
fact $S_3$-symmetric on the three blocks: every permutation in
$S_3$ of the indices $\{1, 2, 3\}$ is an automorphism of
$(\LL, +, \star)$.

\subsection*{Verification}

The two generators of $S_3$ are tested in
\texttt{python\_project/src/probe\_aut\_lambda\_star.py} using the
exact $\ZZ$-basis of $\LL$ from
\texttt{verify\_consecutive\_twists\_exact.py}; for each candidate
$g$ on the 24-vector basis $B$ of $\LL$, the script checks
$g(B_i \star B_j) = g(B_i) \star g(B_j)$ for all $24 \times 24 = 576$
basis pairs.

\medskip
\begin{tabular}{ll}
\toprule
candidate & verdict \\
\midrule
$\ZZ_3$ cyclic block permutation $(x,y,z) \mapsto (y,z,x)$
  & preserves $\LL$ AND $\star$ \\
$S_3$ transposition of blocks $1 \leftrightarrow 2$
  & preserves $\LL$ AND $\star$ \\
componentwise negation $-I_{24}$
  & preserves $\LL$, does NOT preserve $\star$ \\
$\sigma$ applied to each block
  & does NOT preserve $\LL$ \\
\bottomrule
\end{tabular}
\medskip

\noindent
The first two candidates generate $S_3$: a 3-cycle and a
transposition on $\{1, 2, 3\}$ together generate the full $S_3$.
Hence every $S_3$ permutation of the three octonion blocks is in
$\Aut(\LL, +, \star)$.

\subsection*{Why this is more than ``$\ZZ_3$-symmetric''}

The paper's title and the definition of $\star$ both emphasise the
$\ZZ_3$ routing: same multiplication in all three blocks, with a
cyclic $\ZZ_3$ routing of cross-block products.  But $\ZZ_3$ is
only the cyclic group of order three; the empirical finding here
is that the construction is symmetric under the full $S_3$,
which is generated by the cyclic $\ZZ_3$ plus any single
block-transposition.

In other words, the $\ZZ_3$ routing on the right-hand side of
Definition 3.3 is consistent with a stronger $S_3$ symmetry of the
resulting bilinear product on $\RR^{24}$.  This is empirically
visible only by checking a block-transposition explicitly (a
$\ZZ_3$-symmetry argument from the definition alone gives only the
3-cycle).

\subsection*{Bonus: strict subgroup of $\Co_0$}

The same probe also records:
\begin{itemize}
\item $-I_{24} \in \Co_0 \setminus \Aut(\LL, +, \star)$: the
  componentwise negation lies in the Conway group $\Co_0$
  (preserves $\LL$), but does not preserve $\star$ (since
  $\star$ is bilinear, $-I_{24}(u) \star -I_{24}(v) = u \star v$,
  whereas $-I_{24}$ acting on the result of $u \star v$ would give
  $-(u \star v)$; the two disagree unless $u \star v = 0$).
\item Therefore $\Aut(\LL, +, \star) \subsetneq \Co_0$ is a
  \emph{strict} inclusion.
\end{itemize}

This is consistent with the paper's \S7 open question on the
automorphism group of $(\LL, +, \star)$: the question is still
genuinely open, but the bound from below has improved (a copy of
$S_3$ inside $\Aut$ rather than only $\ZZ_3$) and the bound from
above has improved (strict inclusion $\Aut(\LL,+,\star) \subsetneq
\Co_0$).

\subsection*{Proposed handling}

\begin{enumerate}
\item Note the $S_3$ symmetry in v4's \S6 ``\textbf{Comparison.}''
  paragraph, which currently says ``the cross-block $\ZZ_3$
  symmetry is supplied by the routing of Definition 3.3
  rather than by the octonion product itself'' -- replace ``$\ZZ_3$''
  with ``$S_3$'' (or ``$S_3$ extending the $\ZZ_3$ routing'').
\item Alternatively (or in addition), add a one-line remark at the
  end of \S3 noting that the routing in Definition 3.3 induces an
  $S_3$ symmetry on the three blocks, verified empirically.
\item Update \S7's bullet on the automorphism group to record
  $\Aut(\LL, +, \star) \supseteq S_3$ (block permutations) and
  $\Aut(\LL, +, \star) \subsetneq \Co_0$ (strict inclusion from
  $-I_{24}$).
\end{enumerate}

No retitling is needed: the paper's title ``$\ZZ_3$-symmetric
triple-octonion product'' refers to the $\ZZ_3$ routing rule used
to \emph{define} the product, not to the full symmetry of the
resulting object.  $\ZZ_3$ in the title remains accurate as a
description of the construction.

% ===================================================================
\section*{Cross-references}
% ===================================================================

\begin{itemize}
\item Main paper: \texttt{paper/main.tex} (v4, frozen 25 May 2026
  at \texttt{paper/main\_2026-05-25.tex}).
\item v3 $\to$ v4 changelog: \texttt{paper/v3\_to\_v4\_summary.md}.
\item $S_3$ probe script:
  \texttt{python\_project/src/probe\_aut\_lambda\_star.py}.
\item Petersson 2018 (URL preserved in v4 bibliography):
  \texttt{https://www.fernuni-hagen.de/mi/fakultaet/emeriti/docs/petersson/ass.-rem.-int.-oct.pdf}.
\end{itemize}

\end{document}
